\subsection{Αφηρημένο Μεθοδολογικό Μοντέλο και Αρχιτεκτονική Συστήματος}

Η παρούσα ενότητα εισάγει το θεωρητικό υπόβαθρο και την αρχιτεκτονική σύλληψη του προτεινόμενου συστήματος, προσφέροντας μια υψηλού επιπέδου αφαίρεση (abstract level) της μεθοδολογίας μετασχηματισμού νομικών κειμένων σε εκτελέσιμο κώδικα (Rules as Code). Πριν την αναλυτική εξέταση των επιμέρους σταδίων υλοποίησης, κρίνεται απαραίτητη η αποσαφήνιση των θεμελιωδών αρχών λειτουργίας, με έμφαση στη φιλοσοφία του συστήματος, την κατανομή των ρόλων των εμπλεκόμενων οντοτήτων και τις αυστηρές οριακές συνθήκες που διέπουν τη χρήση των Μεγάλων Γλωσσικών Μοντέλων.

\subsubsection{Φιλοσοφία και Κεντρική Ιδέα}
Η βασική σύλληψη του προτεινόμενου συστήματος εδράζεται στη συστηματική και ελεγχόμενη μετάβαση από τον αδόμητο, φυσικό νομικό λόγο σε ένα αυστηρά δομημένο υπολογιστικό περιβάλλον. Σε αντίθεση με θεωρητικές προσεγγίσεις που επιδιώκουν την πλήρη και ανεξέλεγκτη αυτοματοποίηση αυτής της διαδικασίας—συχνά εις βάρος της ακρίβειας ("τυφλή" αυτοματοποίηση)—το παρόν μοντέλο σχεδιάστηκε ως ένα πλαίσιο που θέτει στο απόλυτο επίκεντρο την ιχνηλασιμότητα (traceability). Αποτελεί αδιαπραγμάτευτη αρχή της αρχιτεκτονικής πως κάθε παραγόμενος υπολογιστικός κανόνας, ανεξαρτήτως της πολυπλοκότητάς του, πρέπει να δύναται να συνδεθεί άμεσα, αδιαμφισβήτητα και γραμμικά με το αρχικό νομικό εδάφιο (source text) από το οποίο πηγάζει. Μέσω αυτής της διαρκούς σύνδεσης, διασφαλίζεται η διαφάνεια του συστήματος, παρεμποδίζεται η ενσωμάτωση αυθαίρετων κανόνων και διευκολύνεται καθοριστικά ο έλεγχος νομικής ορθότητας και συμβατότητας.

\subsubsection{Αφηρημένη Ροή Εργασίας (Abstract Workflow)}
Στο πιο αφηρημένο επίπεδο, η μεθοδολογία οπτικοποιείται και λειτουργεί ως ένας αγωγός μετασχηματισμού δεδομένων (data transformation pipeline). Η ροή εργασίας ακολουθεί μια λογική σταδιακής εξειδίκευσης, δόμησης και αποκρυστάλλωσης της διαθέσιμης πληροφορίας. Εκκινεί από το μέγιστο επίπεδο αφαίρεσης, δηλαδή το φυσικό κείμενο (Τμηματοποίηση), και προχωρά στον εντοπισμό και τη σημασιολογική ταξινόμηση των θεμελιωδών συστατικών μερών, όπως τα νομικά υποκείμενα, τα αντικείμενα και οι προϋποθέσεις (Εξαγωγή Οντοτήτων και Κανόνων). Ακολούθως, η αθροισθείσα δομημένη πληροφορία υποβάλλεται σε αυστηρή λογική επαλήθευση (Εντοπισμός Αντιφάσεων). Εφόσον η συνέπεια των κανόνων επιβεβαιωθεί, το σύστημα προχωρά στη μορφοποίησή τους (Μοντελοποίηση Δεδομένων και Παραγωγή Business Rules) και τελικά στη δημιουργία εκτελέσιμης επιχειρησιακής λογικής (Μετασχηματισμός σε GoRules). Κάθε σημείο διεπαφής του αγωγού αφαιρεί περιττή γλωσσική πολυπλοκότητα, προσθέτοντας ταυτόχρονα βαθμούς υπολογιστικής δομής.

\subsubsection{Κατανομή Ρόλων: Το Μοτίβο Human-in-the-Loop (HitL)}
Ακρογωνιαίος λίθος της βιωσιμότητας και της αξιοπιστίας είναι η υιοθέτηση του μοτίβου "Human-in-the-Loop" (HitL). Σε αυτή τη διάταξη, θεσπίζεται ένας αυστηρός διαχωρισμός αρμοδιοτήτων, αποτρέποντας τη λήψη αυτοματοποιημένων κανονιστικών αποφάσεων. Το Μεγάλο Γλωσσικό Μοντέλο (LLM) στοχευμένα αναλαμβάνει τον ρόλο του "γνωσιακού επεξεργαστή" (cognitive engine): είναι επιφορτισμένο με την αναγνώριση περίπλοκων γλωσσικών μοτίβων, την ταχύτατη εξαγωγή πληροφορίας από ογκώδη κείμενα και την ανάδειξη πιθανών λογικών συγκρούσεων.

Αντιθέτως, ο Άνθρωπος —και πιο συγκεκριμένα ο Νομικός Αναλυτής— διατηρεί την αδιαπραγμάτευτη ιδιότητα του επικυρωτή (validator). Καμία μετάβαση δεδομένων από το ένα στάδιο του pipeline στο επόμενο δεν διεκπεραιώνεται αν δεν προηγηθεί συνειδητή ανθρώπινη έγκριση. Ο χρήστης είναι υπεύθυνος να αξιολογήσει τα παραγόμενα αποτελέσματα, να επιλύσει τις αναπόφευκτες ερμηνευτικές ασάφειες του φυσικού λόγου και, τελικά, να προσδώσει στο αποτέλεσμα νομική ισχύ. Το LLM λειτουργεί αμιγώς επικουρικά και ενισχυτικά της ανθρώπινης παραγωγικότητας, δίχως ποτέ να υποκαθιστά τη νομική κρίση.

\subsubsection{Αφηρημένο Σχήμα Εισόδου/Εξόδου (Input/Output Schema Paradigm)}
Από τη σκοπιά της μηχανικής λογισμικού, η αρχιτεκτονική του αγωγού διέπεται από ένα αυστηρό παράδειγμα εισόδου/εξόδου. Η είσοδος (input) του συστήματος συνίσταται πάντοτε σε αδόμητη ή μερικώς δομημένη πληροφορία (unstructured text), υλοποιημένη στα σώματα των νομοθετικών κειμένων. Η εν λόγω αδόμητη δομή οφείλει ωστόσο να συνοδεύεται από κρίσιμα μεταδεδομένα (metadata), όπως η πηγή, το έτος ή το σχετικό άρθρο, τα οποία χρησιμεύουν ως σταθερές συντεταγμένες της πληροφορίας.

Στον αντίποδα, η έξοδος (output) κάθε ενδιάμεσου σταδίου —και θεμελιωδώς η τελική έξοδος— αποτελείται αποκλειστικά από τυποποιημένα, δομημένα δεδομένα αναγνώσιμα από μηχανές (machine-readable schemas), μέσω τεχνολογιών όπως τα αρχεία μορφής JSON. Η θεμελιώδης προδιαγραφή αυτού του σχήματος απαγορεύει τις ορφανές εγγραφές: κάθε εννοιολογικό συστατικό φέρει μοναδικά αναγνωριστικά (IDs) τα οποία το χαρτογραφούν άμεσα στην πηγή του. Η αρχιτεκτονική αυτή εγγυάται πως ό,τι παράγεται, μπορεί άμεσα να ελεγχθεί και να ανακτηθεί από την αρχική του μορφή.

\subsubsection{Οριακές Συνθήκες και Περιορισμοί του LLM (Boundary Conditions)}
Τέλος, καθίσταται απολύτως αναγκαίος ο καθορισμός των οριακών συνθηκών (boundary conditions) λειτουργίας του συστήματος, μέσω της θέσπισης άκαμπτων κανόνων εμπλοκής (guardrails) στο υποκείμενο μοντέλο. Πρώτον, στο LLM απαγορεύεται δραστικά να παράγει αυθαίρετες παραδοχές ή γεγονότα (hallucinations), οφείλοντας να στηρίζεται αποκλειστικά στο παρεχόμενο κείμενο. Δεύτερον, απαγορεύεται συστημικά στο μοντέλο να καλύψει νομοθετικά κενά (lacunae) με δική του, "λογική" πρωτοβουλία: εάν ένας νόμος είναι ελλιπής, η μεταφορά του σε κώδικα πρέπει να αντανακλά την έλλειψη αυτή. 

Η μοναδική αποδεκτή συμπεριφορά του LLM σε περιπτώσεις αντικρουόμενων διατάξεων ή αστοχιών είναι να ενεργεί ως μηχανισμός επισήμανσης (flag). Μέσω συστημικών προειδοποιήσεων ειδοποιεί τον νομικό αναλυτή για την ύπαρξη κάποιας αντίφασης, παρέχοντας στον άνθρωπο-επικυρωτή τα αναγκαία τεκμήρια προκειμένου να λάβει την απόφαση επίλυσης. Κατ' αυτόν τον τρόπο διαφυλάσσεται απολύτως ο επιστημονικός εκδημοκρατισμός και το τελικό κύρος της διαδικασίας παραγωγής των κανόνων.
